\section{Versuchsaufbau}
\label{sec:versuchsaufbau}

Die Evaluation folgt einer gestuften Logik. Ein ressourcenschonender
Smoke-Test prüft zunächst die Datenverträge. Separate COLMAP-Voruntersuchungen
begründen den SfM-Standard. Danach untersuchen automatisierte Matrixläufe
Kameramodell, FPS und Auflösung. Eine kontrollierte Vierfeld-Ablation lokalisiert
den Einfluss von Gaussian-Optimierung und Tiefenroute\footnote{Die
Tiefenroute bezeichnet die Art, wie die für das Surface-Sampling benötigte
Oberflächentiefe gewonnen wird: entweder durch Projektion der
SuGaR-``Diamond''-Primitive in einen Z-Buffer (Diamond-Mesh-Tiefe) oder
direkt aus dem Gaussian-Rasterizer.}. Abschließend
vervollständigt eine zwölf Läufe umfassende SuGaR-Folgeprüfung den technischen
Meshroutenvergleich auf der Coarse-Stufe.

\subsection{Testdatensatz, Variablen und Konstanten}

Die Tests verwenden das Alurohr-Video als linearen Labor-Testdatensatz. Die
Matrixprofile bestehen aus 5 FPS und 2 FPS sowie den Auflösungen 1280$\times$720,
960$\times$540 und 640$\times$360. Der Prompt lautet \texttt{pipe}. SIFT-
Merkmale, Sequential-Overlap, Guided Matching, STS-Iterationszahl, Meshziel,
Surface-Sample-Zahl und Seed werden konstant gehalten.

\begin{table}[H]
\centering
\small
\begin{tabularx}{\textwidth}{p{3.5cm}X}
\toprule
Kategorie & Festlegung\\
\midrule
Datensatz/Prompt & Alurohr-Video, \texttt{pipe}\\
Kameramodelle & \texttt{SIMPLE\_RADIAL}, \texttt{PINHOLE}, \texttt{OPENCV}\\
Frameprofile & 2/5 FPS; $1280\times720$, $960\times540$, $640\times360$\\
COLMAP-Standard & 4096 Plain-SIFT, Overlap 15, Guided Matching aus,
Peak-Threshold 0,003\\
STS & 7000 Gesamtiterationen, 5000 Objektcurriculum\\
Mesh & 200.000 Zielvertices, 5.000.000 Surface-Samples, Poisson-Tiefe 10,
Seed 42\\
Masken & RGB \texttt{default}, DN \texttt{middle}, keine zusätzliche
Dilatation\\
Centerline & \texttt{single}, Voxel 0,1, B-Spline Grad 10\\
Bildmetriken & ausschließlich objektmaskierte PSNR, SSIM und LPIPS\\
\bottomrule
\end{tabularx}
\caption{Zentrale Variablen und kontrollierte Standardparameter.}
\label{tab:versuchsparameter}
\end{table}

Die Bildmetriken sind die einzigen Qualitätsmetriken für Renderings.
Maskenabdeckung, Registrierungsquote, Laufstatus und Laufzeit werden als
technische Zusatzgrößen geführt. Sie werden nicht zu weiteren Bildmetriken
umdefiniert.

Für die Maskenlogik werden zwei Funktionen getrennt: Die erodierte
\texttt{middle}-Maske bildet einen konservativen Objektkern und wird für das
Coverage-Gate sowie für die Auswahl nichtleerer Eval-Ansichten verwendet. Die
eigentliche Bildmetrik wird anschließend mit der weniger beschnittenen
\texttt{default}-Maske berechnet. Dadurch werden unsichere Randbereiche bei
der Zulassung eines Frames streng geprüft, ohne die Bildmetriken unnötig auf
eine erodierte Objektfläche zu beschränken. Das Feld
\texttt{mask\_level=default} in den Metrik-JSONs dokumentiert diese tatsächliche
Auswertungsmaske.

Die Maskenhierarchie selbst ist im Projektverlauf vereinheitlicht worden:
Historisch war zwischenzeitlich auch die \texttt{default}-Maske dilatiert;
verbindlich gilt nun \texttt{default} = unverändert, \texttt{middle} = einfach
erodiert, \texttt{small} = zweifach erodiert, ohne Dilatation. SuGaR nutzt
dabei nicht alle Maskenstufen gleichzeitig: Der maskierte RGB-Verlust der
Coarse-Route arbeitet auf \texttt{default}, die Depth-Normal-Consistency auf
der konservativeren \texttt{middle}-Maske.

\subsection{COLMAP-Voruntersuchung}

Eine getrennte CPU-Voruntersuchung vergleicht die Sparse-SfM-Stufe. Sie umfasst
unter anderem 4096 gegen 8192 SIFT-Merkmale, Guided Matching und zwei
Kompressionsprofile. Alle fünf ausgewerteten Hauptläufe registrierten 240 von
240 Bildern. Die Untersuchung begründet den Standard 4096 Plain-SIFT und
deaktiviertes Guided Matching, ist aber kein vollständiger STS-/Meshvergleich.

\subsection{Smoke-Test und Replay}

Der Lauf \path{matrix_smoke_low_pipe_full} ist ein einzelner Low-
Auflösungstest mit 5 FPS, \texttt{SIMPLE\_RADIAL} und Route A. Er prüft die Kette von den
idealen Masken über einen vorhandenen STS-Checkpoint bis zu Mesh, Rendering,
Metriken und Postprocessing. Der Zusatz \texttt{full} bedeutet, dass dieser
einzelne Lauf nicht nach der Maskenprüfung beendet wurde; er bezeichnet keine
vollständige Parameter-Matrix.

Der Replay startet an einer Schrittgrenze und stellt nur die dafür notwendigen
Artefakte wieder her. Dadurch können teure, bereits validierte Stufen
wiederverwendet werden, ohne einen abgebrochenen Optimierer innerhalb einer
Iteration fortzusetzen.

\subsection{Kameramodell-/FPS-/Auflösungsmatrix}

Die bereits ausgeführte A-Route wurde nicht erneut gerechnet. Die Folge-Matrix
\path{matrix_rest} enthält PINHOLE-A und OPENCV-A. Die gezielte
SuGaR-Folgeprüfung verwendet die separate Konfiguration
\path{tools/experiment_matrix_sugar.tsv} mit SIMPLE\_RADIAL-SuGaR und
OPENCV-SuGaR. Damit werden Kameramodell, FPS und Auflösung getrennt von der
Meshroute untersucht.

Die zusammengeführte Hauptauswertung enthält sechs historische
\texttt{simple\_radial\_a}-Versuche und 18 Folgeexperimente. Die Folgeexperimente
umfassen \texttt{PINHOLE}-A, \texttt{OPENCV}-A und die damalige
\texttt{SIMPLE\_RADIAL}-SuGaR-Route über jeweils drei Auflösungen und zwei
FPS-Profile. Der Runner arbeitet seriell, weil alle Versuche denselben
Live-Workspace verwenden und zwischen den Läufen archiviert und bereinigt wird.

\subsection{Vierfeld-Ablation der Meshgewinnung}

Die Ursache sichtbarer Meshunterschiede wird mit vier Coarse-only-Varianten
getrennt:

\begin{enumerate}
  \item A: Original-STS-GS mit projizierter Diamond-Mesh-Tiefe;
  \item B: Original-STS-GS mit Gaussian-Rasterizer-Tiefe;
  \item C: SuGaR-Coarse $c=9000$ mit Diamond-Mesh-Tiefe;
  \item D: SuGaR-Coarse $c=9000$ mit Gaussian-Rasterizer-Tiefe.
\end{enumerate}

Kameras, Eingabemasken, Surface-Level, Poisson-Tiefe, Dichtequantil,
Vertexziel, Samplezahl, Projektion und Seed bleiben konstant. Dadurch werden
Gaussian-Optimierung und Tiefenroute als getrennte Einflussfaktoren untersucht.
Die aus gleichmäßig gesampelten Meshpunkten berechneten gerichteten Distanzen
sind relative Vergleichsgrößen zwischen den Varianten, keine Abweichungen von
einer Ground Truth.

\subsection{Abgeschlossene SuGaR-Folgeprüfung}

Die Folgematrix umfasst zwölf Läufe über die Kameramodelle
\texttt{SIMPLE\_RADIAL} und \texttt{OPENCV} (jeweils mit
\texttt{sugar\_coarse}) bei drei Auflösungen und zwei FPS-Profilen. Alle zwölf
Manifeste und Coarse-Metrikdateien sind erfolgreich
archiviert. Ein verfeinertes PLY wurde in diesen Läufen nicht exportiert; die
Folgematrix bewertet daher nur die Coarse-Route. Die Ergebnisse zeigen die
Grafiken in Abschnitt~\ref{sec:sugar-folgepruefung} und die Fehlerchronik in
Tabelle~\ref{tab:fehlerchronik}.

\subsection{Evaluationsprotokoll}

Für jeden Lauf werden Masken- und Warp-Coverage sowie die COLMAP-Ausgaben
archiviert. Hinzu kommen STS-Checkpoint (\path{point_cloud.ply}), gegebenenfalls
SuGaR-Coarse-Checkpoint (\path{9000.pt}), Mesh-/Postprocess-Ausgaben, feste
Eval-Frames und objektmaskierte PSNR-, SSIM- und LPIPS-Werte. Die Daten werden über
\texttt{manifest.json}, \texttt{parameters.json}, \texttt{run.log} und
\texttt{run.md} nachvollziehbar gehalten.

Die Matrixgrafiken verwenden den arithmetischen Mittelwert aus 2 und 5 FPS nur,
wenn beide Läufe derselben Konfiguration vollständig erfolgreich waren. Ein
technischer Fehler wird nicht als schlechter Qualitätswert in einen Mittelwert
umgewandelt. Die neue Auswertung trennt zusätzlich die Modellstufen: Eine
STS-Grafik liest ausschließlich die aus \path{point_cloud.ply} gerenderten
\texttt{sts\_masked.json}-Dateien (Route~A), eine SuGaR-Grafik ausschließlich die
aus \path{9000.pt} gerenderten \texttt{sugar\_coarse\_masked.json}-Dateien. Da
beide Formate dieselben 3D-Gaussian-Parameter speichern, beruht der Unterschied
in den Messwerten (rund $29{,}5\,\mathrm{dB}$ gegenüber $21{,}5\,\mathrm{dB}$)
rein auf den veränderten Gaussian-Parametern nach der SuGaR-Coarse-Optimierung und
nicht auf dem Dateiformat. Die \texttt{per\_frame}-Werte werden ergänzend als
Boxplots mit Einzelpunkten dargestellt, damit Streuung und Ausreißer nicht durch
den Mittelwert verborgen werden.

Ein Lauf erhält nur dann den Status vollständig erfolgreich, wenn alle für
seine Route vorgesehenen Stufen gültig sind: Masken und Coverage,
COLMAP-Modell, Roh-/Ideal-Domäne, vollständig gematchter Eval-Split,
STS-Checkpoint, Objektselektion, Mesh, Postprocessing, objektmaskierte Metriken
und Ergebnisarchiv. Für SuGaR-Coarse muss zusätzlich das SuGaR-Rendering
erfolgreich sein. Eine zuvor erzeugte \texttt{sts\_masked.json} bleibt eine
STS-Baseline und wird nicht als SuGaR-Messung gewertet.

\subsection{Erfolgskriterien der Projektarbeit}

Die PA gilt als technisch erfolgreich, wenn:

\begin{enumerate}
  \item die zentrale Pipeline für den Testdatensatz durchläuft;
  \item die Bild- und Maskendomänen konsistent sind;
  \item STS, Meshgewinnung und Centerline-Extraktion reproduzierbar arbeiten;
  \item der Status jedes Matrixlaufs und die Fehlerursache archiviert sind;
  \item die Ergebnisse für eine spätere GNSS-basierte Bachelorbewertung
        geeignet abgelegt werden.
    \item interaktiver Vollauf, Replay, Autopilot und Matrixrunner in ihrer
      jeweiligen Funktion eindeutig dokumentiert und mindestens über einen
      archivierten Lauf oder Vertragstest nachgewiesen sind.
\end{enumerate}

Die PSNR-/SSIM-/LPIPS-Werte sind dabei Funktions- und Visualisierungsmetriken.
Sie ersetzen nicht die spätere geometrische Bewertung mit Centerline-RMSE,
Hausdorff-Distanz und GNSS-Referenz.
